\documentclass[11pt,letterpaper]{article}
\usepackage[margin=0.7in]{geometry}
\usepackage[T1]{fontenc}
\usepackage{lmodern,amsmath,amssymb,graphicx,array,booktabs,caption,pdflscape,float,microtype,xcolor,hyperref}
\hypersetup{colorlinks=true,linkcolor=blue!50!black,urlcolor=blue!50!black,citecolor=blue!50!black,
pdftitle={HPS v5.7.1: Magnetic transport and dark-photon acceptance}}
\setlength{\parskip}{5pt}
\setlength{\parindent}{0pt}
\captionsetup{font=small,labelfont=bf}
\newcommand{\vFiveSevenRoot}{.}
\begin{document}
{\Large\bfseries Magnetic transport and dark-photon acceptance in HPS}\par
{\large HPS-GPR v5.7.1}\par
13 September 2026

The opening angles of an $A'\to e^+e^-$ decay decrease as its laboratory
energy increases. This shifts the mass scale at which the daughter tracks
enter a forward spectrometer. The calculation below connects that boost to
the active silicon geometry and magnetic fields for representative 2015,
2016 and 2021 HPS configurations. The electron and positron are propagated
through the field maps and tested against finite sensor faces using the
run-dependent hit requirements. This replaces the assumed 70 mrad outer
angle of v5.7.0 with a calculation of curved tracks through the detector.

The resulting fractions describe prompt decays at a fixed parent energy and
direction. They include magnetic transport and the hit component of the
selection; material interactions, hit finding, trigger, reconstruction and
other analysis cuts remain outside the calculation. Zero-field curves are
retained to show the effect of bending. The equations and detector assumptions
precede the figures so this section can be inserted into a longer analysis note.

% Insertable body; set \vFiveSevenRoot to the copied source directory.
% Required packages: amsmath, amssymb, graphicx, booktabs, caption, pdflscape, float.
\providecommand{\vFiveSevenRoot}{.}
\input{\vFiveSevenRoot/sections/kinematics.tex}
\input{\vFiveSevenRoot/sections/geometry_results.tex}

\clearpage
\begin{thebibliography}{99}
\bibitem{v571:HPS2022} N. Baltzell et al., \emph{The Heavy Photon Search Experiment},
\href{https://arxiv.org/abs/2203.08324}{arXiv:2203.08324} (2022), Table I and detector description.
\bibitem{v571:HPS2015} P. H. Adrian et al., \emph{Search for a dark photon in electroproduced
$e^+e^-$ pairs with HPS}, Phys. Rev. D \textbf{98}, 091101 (2018),
\href{https://arxiv.org/abs/1807.11530}{arXiv:1807.11530}.
\bibitem{v571:HPS2016} P. H. Adrian et al., \emph{Searching for prompt and long-lived dark
photons with HPS}, Phys. Rev. D \textbf{108}, 012015 (2023),
\href{https://arxiv.org/abs/2212.10629}{arXiv:2212.10629}.
\bibitem{v571:geometry-source} Jefferson Lab, \emph{hps-java detector descriptions},
revision \texttt{473c732d},
\href{https://github.com/JeffersonLab/hps-java/tree/473c732dafaebdd3f58389cd1052a35c60098a23/detector-data/detectors}{detector-data/detectors}.
Committed compact and LCDD snapshots and their SHA-256 hashes are included in this package.
\bibitem{v571:beam-coords} Jefferson Lab, \emph{hps-mc beam coordinate conversion},
revision \texttt{308e27be},
\href{https://github.com/JeffersonLab/hps-mc/blob/308e27be7821f1aad406cb8060602a10fca2d37a/tools/stdhep-tools/src/beam_coords.cc}{beam\_coords.cc}.
\bibitem{v571:test-detector} M. Battaglieri et al., \emph{The Heavy Photon Search Test Detector},
\href{https://arxiv.org/abs/1406.6115}{arXiv:1406.6115}.
The historical approximate 15--70 mrad design interval motivated the original v5.7.0 illustration;
it is not used as the finite-geometry boundary in v5.7.1.
\bibitem{v571:internal2015} O. Moreno, N. Baltzell, M. Graham and J. Jaros,
\emph{Search for a Heavy Photon in Electro-Produced $e^+e^-$ Pairs with the HPS Experiment at JLab},
archived HPS internal note (undated local copy), PDF pp.~11--12, Table~3.
The local file hash and source locators are recorded in
\nolinkurl{inputs/selection/selection_criteria.json}.
\bibitem{v571:note504} E. Peets, M. Gignac and E. Hsu,
\emph{HPS Gaussian-Process Resonance Search: Analysis Note v5.0.4},
local analysis-note snapshot dated 11 September 2026, event-selection section and table.
Unmodified source excerpts are included in \nolinkurl{inputs/selection/};
the current user correction supersedes their older 2021 hit-count statements for this study.
\bibitem{v571:fieldmaps} Jefferson Lab, \emph{HPS magnetic-field maps},
revision \texttt{feadfcbd}, \href{https://github.com/JeffersonLab/hps-fieldmaps/blob/feadfcbd17bdee8515cd962d75a90ba1534fa888/README.md}{run-period catalogue}.
Active archives: \href{https://raw.githubusercontent.com/JeffersonLab/hps-fieldmaps/feadfcbd17bdee8515cd962d75a90ba1534fa888/125acm2_3kg_corrected_unfolded_scaled_0.7992_v3.tar.gz}{2015 corrected v3},
\href{https://raw.githubusercontent.com/JeffersonLab/hps-fieldmaps/feadfcbd17bdee8515cd962d75a90ba1534fa888/209acm2_5kg_corrected_unfolded_scaled_1.04545_v4.tar.gz}{2016 v4},
\href{https://raw.githubusercontent.com/JeffersonLab/hps-fieldmaps/feadfcbd17bdee8515cd962d75a90ba1534fa888/334acm3_8kg_corrected_unfolded_scaled_1.07326.tar.gz}{2021 data map}.
Archive and extracted-file hashes are recorded with the active field configuration.
\bibitem{v571:fieldreaders} Jefferson Lab, \emph{hps-lcsim field-map reader},
revision \texttt{391d4649}, \href{https://github.com/JeffersonLab/hps-lcsim/blob/391d46492257c5ab155e22ecf74734efd7aaa27a/detector-framework/src/main/java/org/lcsim/geometry/field/FieldMap3D.java}{FieldMap3D.java};
SLAC, \emph{LCDD Cartesian field-map reader}, revision \texttt{36a4805e},
\href{https://github.com/slaclab/lcdd/blob/36a4805e2ff82ec04f7e5332532ef751fddb5057/src/lcdd/bfield/Cartesian3DMagneticFieldMap.cc}{Cartesian3DMagneticFieldMap.cc}.
These sources establish field units, global-offset subtraction and interpolation.
\end{thebibliography}

\end{document}
